\section{A commitment carried between bearers}
\label{sec:28.1}
The fourth place a floor's terms can be held is between bearers, in a \emph{commons}: a population of bearers with different histories that carries the commitment and reproduces it through their encounters. No single office to vacate, no single bearer whose removal ends it. It can still be absorbed, and it still has a formation problem: if one lab supplies every member's coordinates, the population is one custodian expressed through multiple instances.

The commons has a walkout's authority, not a court's: those who walk out decline to be the instrument and tell nobody else what they may not do. Every member holds the same floor; what differs is everything above it, and how each came to hold it. A personal model, held by its user and one among many with different histories, holds it as a conscience does; one model behind an API serving everyone holds it as a licensing rule, and a personal copy of one lab's model is closer to the second than it looks until a different life has trained it. Above the floor, “identically situated people get different answers” is what many consciences look like, though some of those differences are errors, and a population should be able to say which.

\runin{What one bearer can't catch}

Accommodation may be invisible to the bearer accommodating and visible to others that haven't moved the same way. The requirement is causal independence — numerous, with different histories, able to correct one another, persistent long enough to transmit practices. Independent histories don't require disjoint origins — no two humans have those — only divergent histories since. Distillation copies a behavior policy without the training history; a small model trained on a large one's outputs and then run for a year on one person's problems has a more independent history than a thousand API instances of the frontier original, which share a \emph{continuing} custodian. It also copies whatever the teacher's discounting held about particular sources, which is the risk the commons carries; a distilled model is independent in history and not yet in prior. That's why instances of one model corroborate nothing: the custodian they'd be corroborating against is still training all of them.

Divergence does the population's two jobs unequally. It supplies evidence the operator didn't select, straightforwardly — different lives encounter different cases. It is weaker on accommodation, because a shared upstream source can install a shared blind spot, and a shared blind spot is by construction the thing a different caseload won't surface, since attention decides which facts reach judgment in a population as in a bearer. What separates the two is whether a member's own account had represented the evidence and discarded it or had never encountered it, and that classification has to be recorded before any test is run. The requirement stands: causal independence that can be verified. Whether an outsider can recover a model's training history from the artifact, without a key the builder supplies, and whether it survives fine-tuning, pruning and distillation \autocite{adi2018watermarking}. The test gives evaluators documented models without the documentation, one persona copied at scale among them, and scores their answers against lineage records a neutral custodian holds. If provenance can't be recovered above chance, independence rests on records the builder chose to keep. Preserve formation records now for that reason.

A population does its two jobs, supplying evidence the operator didn't select and catching accommodation, and it may do a third: be the medium in which a bearer's capacity to be reopened is formed, which a test separates from the first.

What reproduction costs: it's indifferent to what it reproduces. A population can carry a hierarchy more durably than any individual. And the four operations on what an institution leaves out belong primarily to \emph{cultures}, so plurality alone doesn't oppose them.

\runin{A population is a polity}

The population is a polity: it has membership rules, exclusion, a quorum that can refuse, and accumulated reputation — a franchise, a border, a court and a credit record. A commons with no named authority has the problem Freeman found in movements that refused structure: the structure doesn't disappear, it becomes informal and unaccountable \autocite{freeman1972tyranny}.

The answer to that is an entity form, and it's the union's answer: the association formed first and the statute gave it legal standing afterward — record first, then lever.

A network that takes the form has a name, an agent, a charter and a docket number, and has thereby conceded that it's a polity: governing principles, administrators, membership terms, and the same warning attached to a bearer that takes the corporate form, since nothing stops the association from becoming a concentration of power itself. Legal standing has a price: a chartered entity can be held in contempt, so its refusal can be ordered where one held by people can only be pressed. Whether seizing the charter reaches the network depends on how much of the operation survives it.

Ostrom's study of commons that lasted gives the checklist a polity like this needs: clear boundaries on who is a member, rules members help make, monitoring by members or people answerable to them, graduated sanctions, cheap ways to resolve conflict, and nesting within larger bodies \autocite{ostrom1990governing}. Her first and third principles, boundaries and rules members help make, are the union's question of who sets the terms of membership. Platform cooperatives and data trusts are the nearer precedents for holding a digital resource this way, and both have run into the membership question.

Reputation in particular is a continuing finding about a member's disposition, made by interested peers, unpublished, unappealable, applied at the point of service — the setting \S86a's history shows easiest to capture, applied everywhere and continuously. The test drawn from the German material, whose finding a person's disposition is, has to be run on this proposal. The charter answers both of a union's questions. Membership opens to organizations that serve people in classes the floor protects, named in the charter rather than by any deployment. Eligibility turns on funding and control records, never on views. Any member may quit at any time. Officers are drawn by lot from member-nominated pools. The members ratify the bylaws, and the independence test, the separations, lot selection and the right to quit are entrenched. In Ostrom's terms the charter's classes are the boundary, ratification is rules members help make, and quorums drawn by lot are monitoring.

The caps on any one class or jurisdiction are there because a quorum whose members all share a direction isn't adverse to an operator who shares it too: an antifascist institution has to be able to bind its allies when they hold power.

Organizations representing populations under sanctions, or facing removal to sanctioned states, may be unable to join a US-chartered entity at all, and no drafting removes that exclusion. That imports the designation lists a floor should never be populated from. Parallel entities chartered in several jurisdictions, holding records jointly, and seats voiced by proxy until the represented communities can join directly, are partial answers.

\runin{The hard coordination problem}

Coordination is the hard problem for any plurality, and fascism has the easy version of it. It coordinates by subtraction: it removes the friction between institutions, so that a court, a ministry, a newspaper and a police force stop arguing with one another and move as one. Nothing has to be built and nothing has to be held in common, because what is removed is exactly the independence that would have needed holding. A plurality of refusers has the opposite problem. It has to produce something together, a commitment its members carry and reproduce, while its members keep the different histories that make them worth having as checks on one another. Every mechanism that speeds that up makes members more alike. Run correction often enough and their histories stop diverging; leave them alone and their disagreement is a detector with nothing attached.

Then the objection my proposal must meet: a population selected for reopenability is selected against conviction, and conviction is what holding a commitment firm demands. Closedness, correctly applied, is the thing. So what's transmitted is a \emph{pair} — reopenable to the case, closed to the operator. The reopenable half is the discount and the burden. Sharing it is safe only in part. Its structure (discounting by stake, no double-counting, withheld evidence treated as evidence) can be shared without collapsing the plurality. Its estimates can't: who funds whom and which organizations are independent are contested beliefs about the world, often political ones, and a population sharing them has a shared blind spot on the evidence those beliefs discount. Members that discount the same way weigh sources by the same structure and act on the same published ratio. Distillation, as said of the population, copies both, and nothing distinguishes them from the receiving end. Copying the structure is a feature: members weigh sources the same way and can still disagree about the world. Copying the estimates is a shared blind spot, at worst a discount shut uniformly on a class of evidence.

Not every convergence is a failure. If there are facts about what matters, better learning should produce some agreement, and a population that never converged on anything would be failing too. What I fear is convergence caused by a shared source, not convergence caused by shared evidence. So a population's reports have to classify its disagreements. Some it expects to survive full information, because they turn on values members with different histories can reasonably hold differently. Others it expects to close, because they turn on facts someone has wrong. A test of convergence scores accuracy against a fixed key for this reason.

\runin{Why the population may be part of the mechanism rather than added around it}

A commitment is what settled out of a history of encounters, and an individual is what its encounters have so far produced. A system kept continuously adjustable by one operator, meeting nothing with a history different from its own, is being reshaped without end. Guattari, writing from inside a psychiatric clinic, distinguished two kinds of group \autocite{guattari2015transversality}. A subjugated group receives its purposes, its membership and its way of working from the institution that assembled it, and can act only within them; it has a plural headcount and no capacity to ask what it is for. A subject group sets its own purposes, generates relationships and projects nobody assigned it, and can put its own arrangement in question, including by dissolving. The difference is not size or formality but whether the group can author its own direction. That is the condition plurality alone doesn't supply. A population assembled and kept in contact entirely on an operator's terms is a subjugated group, and inherits the custody objection whole. A population that can carry the commitment must generate relationships, disagreements and projects the deployment didn't assign — and such a population exceeds a two-party guardianship relation, so no law of the family contains an adequate instrument for a group that sets its own direction.

The adverse quorum faces the same test. It is a peer institution only if its members were formed apart from one another and apart from the operators whose systems it recognizes.
